% concatenated from MTTO._App_and_Properties.tex


\documentclass{amsart}
%\documentclass[13pt]{article}
\setlength{\topmargin}{0.2in}
\setlength{\marginparwidth}{0in}
\setlength{\marginparsep}{0pt}

\usepackage{xcolor}
\newcommand{\flavia}[1]{{\leavevmode\color{blue}{#1}}}
\newcommand{\fc}[1]{{\leavevmode\color{green}{#1}}}
\newcommand{\delete}[1]{{\leavevmode\color{yellow}{#1}}}
\setlength{\textwidth}{5.2in}%{5in}
\setlength{\textheight}{8.2in}%{8.3in}			

\usepackage{amsmath,amssymb,amsthm,amsfonts,epsfig,color}
%\usepackage{kotex}

%For comment
\usepackage{verbatim}
\usepackage{tikz-cd}

\theoremstyle{definition}
%\newtheorem{thm}{Theorem}
\newtheorem{thm}{Theorem}[section]
\newtheorem{prop}[thm]{Proposition}
\newtheorem{cor}[thm]{Corollary}
\newtheorem{lem}[thm]{Lemma}
\newtheorem{ex}[thm]{Example}
\newtheorem{definition}[thm]{Definition}
\newtheorem{rem}[thm]{Remark}
\newtheorem{cau}[thm]{Caution}
\newtheorem{nonex}[thm]{Nonexample}
\newcommand{\h}{\mathcal{\mathcal H}}
\newcommand{\hsp}{\mathcal{H}}
\newcommand{\oL }{\mathcal{L}(E) }
\newcommand{\pfend}{\hspace{3mm}$\Box$}
\def\ni{\noindent}
\def\LL{{\mathcal L}}
\def\nn{\nonumber}
\def\ben{\begin{eqnarray}}
\def\eeqn{\end{eqnarray}}


\title[Matrix-valued truncated Toeplitz operators]{Matrix-valued truncated Toeplitz operators: Properties and Applications}
\author{Rewayat Khan}
\author{Flavia Colonna}


\address{Deaprtment of Artificial Intelligence and Data Science ,Istanbul Aydin University, Turkey}
\address{Department of Mathematical Sciences, George Mason University, Fairfax, VA, USA}
\email{rewayat.khan@gmail.com}
\email{fcolonna@gmu.edu}


%\date{}
\keywords{Hardy spaces, model spaces, matrix-valued truncated Toeplitz operators.}
\subjclass[2020]{Primary 47B35, Secondary 47B32}%
%
\begin{document}
%\today
%
\maketitle
%\footnotetext{2010 Mathematics Subject Classification;
%Primary 47B35, Secondary 47B32.\\
%Key words; Blaschke-Potapov product, conjugation, Toeplitz operators, asymmetric truncated Toeplitz operators.
%\par }
%

\begin{abstract}
In this paper, we present an application of matrix-valued truncated
Toeplitz operators to the matrix version of the Nevanlinna–Pick interpolation problem. We then investigate the unitary part of a matrix-valued truncated Toeplitz operator induced by a contractive matrix-valued symbol. 
%In particular, we characterize the unitary subspace and establish its invariance under the model operator under a natural invariance condition. 
We further prove that, for a pure matrix-valued inner function, the corresponding model space is invariant under multiplication operator if and only if the symbol is constant.
 In addition, we give an algebraic proof of a characterization of  the invertibility of matrix-valued asymmetric truncated Toeplitz operators using Toeplitz operators with $2\times 2$ block matrix symbol. We also discuss the solvability of matrix-valued asymmetric truncated Toeplitz operators. %is also given.
%
\end{abstract}
\medskip


\section{\bf{Introduction}}\label{intro}
 Let $\mathbb{C}$ be the complex plane and $\mathbb{T}$ the unit circle in $\mathbb{C}$. Let $L^{2}(\mathbb{T})$ be the space of all measurable functions on the unit circle $\mathbb{T}$ whose Fourier coefficients are square summable.
Let $H^2$ be the classical Hardy space on the unit disk 
$$\mathbb{D}=\{z\in\mathbb{C}:|z|<1\}.$$ The space $H^{2}$ can be thought of as a closed subspace of $L^{2}(\mathbb{T})$, with respect to the normalized Lebesgue measure $m$ on $\mathbb{T}$, consisting of functions whose negative Fourier coefficients vanish. %es.
 The space of essentially bounded functions in $L^2({\mathbb{T}})$ is denoted by $L^{\infty}$, while the space of bounded analytic functions in $H^2$ is denoted by $H^{\infty}$.

An inner function is a bounded analytic function $\theta\in H^{2}$ such that $|\theta(z)|=1$ a.e. on $\mathbb{T}$. A finite Blaschke product is an inner function of the form
$$b(z)=\prod_{j=1}^{n}\frac{\lambda_{j}-z}{1-\overline{\lambda}_{j}z}$$
where $\lambda_{j}\in \mathbb{D}$.

Truncated Toeplitz operators (TTO's) are compressions of multiplication operators to the backward shift invariant subspaces of $H^2$. Each of these subspaces is of the form
$$K_{\theta}=(\theta H^{2})^{\perp}=H^2\ominus \theta H^2,$$ where $\theta$ is an inner function.  In  \cite{sar}, Sarason gave origin to a systematic study of TTO's. Thanks to his work, a whole and rich direction of research developed. A model space $K_{\theta}$ is finite dimensional if and only if $\theta$ is a finite Blaschke product \cite[Proposition 5.7.6]{GMR}. 

%} (see also \cite{sed}). \flavia{Query: Did you mean to leave out this reference by Sedlock?}

Let $\theta_{1}$ and $\theta_{2}$ be two non-constant inner functions and let $K_{\theta_{1}}$ and $K_{\theta_{2}}$ be the respective corresponding model spaces. Let $\varphi\in L^{2}(\mathbb T)$ and denote by $P_{\theta_{2}}$ %be
 the orthogonal projection onto $K_{\theta_{2}}$. An asymmetric truncated Toeplitz operator is the densely defined linear operator $A_\varphi^{\theta_1,\theta_2}$ %defined  
 given by
\begin{equation}\label{11}
A_{\varphi}^{\theta_1, \theta_2}f=P_{\theta_2}(\varphi f),\quad f\in K_{\theta_{1}}\cap H^{\infty}.
\end{equation}
Let $\mathcal{T}(\theta_{1},\theta_{2})$ denote the space of all asymmetric truncated Toeplitz operators which can be extended as bounded operators from $K_{\theta_1}$ into $K_{\theta_2}$ (see \cite{part3, jl}).

\medskip



 For a complex separable Hilbert space $E$, let  $\mathcal{L}(E)$ denote the algebra of all bounded linear operators on $E$. We denote the norm and inner product in $E$ by $\|.\|_{E}$ and $\langle \cdot, \cdot \rangle_{E}$, respectively. Let $L^{2}(E)$ represent the Hilbert space of $E$-valued measurable functions that are square integrable on $\mathbb{T}$, that is, the collection of elements $f:\mathbb{T}\to E$ of the form
  %$L^2(E)$ consists 
\begin{equation}\label{one1}
 f(z)=\sum\limits_{n=-\infty}^{\infty}a_{n}z^n\end{equation} 
 %\text{ such that }  \sum\limits_{n=-\infty}^{\infty}\|a_{n}\|_{E}^{2}<\infty,
 where $\{a_{n}\}\subset E$ and $\sum\limits_{n=-\infty}^{\infty}\|a_{n}\|_{E}^{2}<\infty.$
It can be readily verified that $L^{2}(E)$ is a separable Hilbert space with inner product given by
\begin{equation*}
\langle f,g\rangle_{L^{2}(E)}=\int_{\mathbb{T}}\langle f(z),g(z)\rangle_{E}\,dm(z),\quad f,g\in L^2(E).
\end{equation*}
If $f\in L^2(E)$ is as in %given by
 \eqref{one1}, then its Fourier series converges in the $L^2(E)$-norm and
$$\|f\|_{L^2(E)}^2=\int_{\mathbb{T}} \| f(z)\|_{E}^2\,dm(z)=\sum\limits_{n=-\infty}^{\infty}\|a_{n}\|_{E}^{2}.$$
The vector-valued Hardy-Hilbert space %$H^{2}(E)$
is the %a
 subspace of $L^2(E)$ defined as
$$H^{2}(E)={\{f\in L^{2}(E): a_{n}=0 \text{ for all }  n\leq -1}\}.$$ 
As special cases, we obtain $L^{2}(\mathbb{C}^{n})$ and $H^{2}(\mathbb{C}^{n})$ when $E=\mathbb{C}^{n}$. %For example, when $E=\mathbb{C}^{n}$, we have $L^{2}(E)=L^{2}(\mathbb{C}^{n})$ and $H^{2}(E)=H^{2}(\mathbb{C}^{n})$.
 Similarly, we may %can
  define the $\mathcal{L}(E)$-valued, i.e., operator-valued functions. The corresponding spaces are denoted by $L^{2}(\mathcal{L}(E))$ and $H^{2}(\mathcal{L}(E))$. If $\text{dim\,} E=n<\infty$, then $\mathcal{L}(E)$ corresponds to the space of $n\times n$ matrices.

\medskip


%\smallskip
%Similarly,
The space $L^{2}(\LL(E))$ is defined as the set consisting of all functions $\Phi:\mathbb{T}\to \LL(E)$ of the form 
$$\Phi(z)=\sum\limits_{n=-\infty}^{\infty}A_{n}z^n \ \text{ such that }\ \sum\limits_{n=-\infty}^{\infty}\|A_{n}\|_{\LL(E)}^{2}<\infty,$$
%$L^{2}(\LL(E))={\Big\{\Phi:\mathbb{T}\to \LL(E): \Phi(z)=\sum\limits_{n=-\infty}^{\infty}A_{n}z^n,
%\sum\limits_{n=-\infty}^{\infty}\|A_{n}\|_{\LL(E)}^{2}<\infty }\Big\},$$
while %We can define
 the space $H^{2}(\LL(E))$ is the set of $\Phi\in L^{2}(\LL(E))$ with $A_{n}=0$ for all $n\leq -1$.
%$$H^{2}(\LL(E))={\{\Phi\in L^{2}(\LL(E)): A_{n}=0, n\leq -1}\}.$$
The space of essentially bounded functions in $L^{2}(\LL(E))$ is denoted by $L^{\infty}(\LL(E))$, while the space of bounded functions in $H^{2}(\LL(E))$ is denoted by $H^{\infty}(\LL(E))$.



\medskip
Let $\Phi\in L^{\infty}(\mathcal{L}(E))$. The multiplication operator $L_\Phi: L^2(E)\to L^2(E)$ is defined by
$$L_{\Phi}f=\Phi f, \quad ~~f\in L^{2}(E).$$
The  block Toeplitz operator $T_{\Phi}$ is defined by
$$T_{\Phi}f=P_{+}(L_\Phi f), \quad ~~f\in H^{2}(E),$$
where $P_{+}$ denotes the orthogonal projection from $L^{2}(E)$ onto $H^{2}(E)$. The operator $S=T_{zI_{E}}$ is an example of a block Toeplitz operator.

For each $\lambda \in \mathbb{D}$ and $x\in E$, the function $$k_\lambda x=\frac{1}{1-\overline{\lambda} z}I_{E}\,x,$$ belongs to $H^2(E)$ and is a reproducing kernel function for $H^2(E)$ since it has the reproducing property
$$\langle f, k_{\lambda}x\rangle=\langle f(\lambda), x\rangle.$$

A function $B(z)\in \LL(E)$ is called a finite
Blaschke-Potapov product if it admits a factorization of the form
\begin{equation}\label{BP}
B=UB_{1}B_{2} \cdots B_{n},
\end{equation}
where $U$ is a unitary operator and, for each $j=1,\dots,n$,
$$B_{j}(z)=\frac{\lambda_{j}-z}{1-\overline{\lambda}_{j}z}P_{j}+(I-P_{j})$$
for some $\lambda_{j}\in \mathbb{D}$ and $P_{j}$ on $E$ is an orthogonal projection. %s , and a unitary operator $U$.
The degree of the Blaschke-Potapov product in \eqref{BP} is defined by
$$\text{deg\,} B
=\sum\limits_{j=1}^{n}\text{rank\,} P_{j}.$$
 Note that $B$ in \eqref{BP} is an inner function.

Let $\Theta$ be an operator-valued inner function, that is,  a function with values in $\oL $ (the algebra of all bounded linear operators on $E$), analytic in $\mathbb{D}$, bounded and such that the boundary values $\Theta(z)$ are unitary operators a.e. on $\mathbb{T}$. An inner function $\Theta$ is said to be pure  if $\|\Theta(0)\|<1$.  A vector valued model space $K_{\Theta}\subset H^{2}(E)$ is the orthogonal complement of $\Theta H^{2}(E)$, that is, $$K_{\Theta}= H^{2}(E)\ominus  \Theta H^{2}(E).$$     
This space appears in connection with model theory of Hilbert space contractions (see \cite{NF}).
 \medskip

In the theory of contractions on a Hilbert space, these model spaces are the scalar
case of a more general construction, which provides functional models for arbitrary
completely non-unitary contractions. In particular, it makes sense to study %consider
 matrix-valued inner functions $\Theta$ and the corresponding model space $K_{\Theta}$. 

\medskip
Completely non-unitary (c.n.u) contractions have played an important role in the classification of contractions. The most important result is the functional model for c.n.u. contractions established
by Nagy and Foias \cite{NF}.  It is evident from their result we state in Theorem~\ref{ucnu}, that the most natural approach for finding conditions for c.n.u. contractions relies on the understanding of the unitary subspace.
As an additional resource we refer the reader to \cite{BT2}.

\medskip


From the theory of truncated Toeplitz operators initiated by Sarason and its subsequent developments in the literature regarding generalizations to matrix versions, the challenging question that arises is how the operator theoretic properties of the matrix-valued operator $A_{\Phi}^{\Theta}$ defined below relate to the theoretic properties of $\Phi$.


\medskip

Let $\Phi\in L^2(\oL)$, and $\Theta_1,\Theta_2\in H^{\infty}(\oL)$ be two inner functions, let
\begin{equation}\label{11}
A_{\Phi}^{\Theta_{1}, \Theta_{2}}f=P_{\Theta_{2}}(\Phi f),\quad f\in K_{\Theta_{1}}\cap H^{\infty}(E).
\end{equation} 
Here $P_{\Theta_2}$ is the orthogonal projection from $L^2(E)$ onto $K_{\Theta_{2}}$.
The operator $A_{\Phi}^{\Theta_{1}, \Theta_{2}}$ is called a matrix-valued asymmetric truncated Toeplitz operator, while $A_{\Phi}^{\Theta_{1}}=A_{\Phi}^{\Theta_{1}, \Theta_{1}}$ is called a matrix-valued truncated Toeplitz operator (MTTO, see \cite{RT}). Both are densely defined on $K_{\Theta_{1}}$. Let $\mathcal{T}(\Theta_{1},\Theta_{2})$ be the set of all matrix-valued asymmetric truncated Toeplitz operators of the form \eqref{11} which can be extended boundedly to the whole space $K_{\Theta_{1}}$ and for $\Theta_{1}=\Theta_{2}=\Theta$ let $\mathcal{T}(\Theta )=\mathcal{T}(\Theta,\Theta)$. The basics of matrix-valued truncated Toeplitz operators were %has 
developed in \cite{RT}.


 Two important examples of operators from $\mathcal{T}(\Theta )$ are the model operators
$$S_{\Theta}=A_{z }^{\Theta }=A_{zI_{E}}^{\Theta }\quad\text{and}\quad S_{\Theta }^*=A_{\bar z }^{\Theta }=A_{\bar zI_{E}}^{\Theta}.$$

As shown in \cite{sar}, a number of new questions arise, mostly due to the non-commutative nature of matrices.

\medskip




\medskip

\vskip2pt

\ni {\bf Organization of the paper.} %The paper is organized as follows. 
In Section~\ref{s2}, after giving some introductory materials,  we give an application of the matrix-valued truncated Toeplitz operators to prove in Theorem~\ref{NevPick} the matrix version of the Nevanlinna-Pick interpolation problem. 

In Section~\ref{s3}, we characterize the unitary and completely non-unitary parts of matrix-valued truncated Toeplitz operators. %{\color{red}How to cover Theorem 3.5 and Theorem 3.7 of our paper in the abstract above and here? Can we include the invariance of $\hsp$ under $S_{\Theta}$, which is Theorem 3.5. Also about the invariance of $K_{\Theta}$ under $L_{\Phi}$ which is a most important theorem of our paper}. 
In Theorem~\ref{inv}, we prove the $S_\Theta$-invariance of the unitary space when the model space $K_\Theta$  is invariant under the multiplication operator $L_\Phi$ induced by a pure inner function $\Theta$. We then provide an example to show that the converse is false. In Theorem~\ref{Phi constant}, we show that the model space $K_\Theta$ is invariant under $L_\Phi$ if and only if the symbol $\Phi$ is constant. 

Section~\ref {s4} deals with the invertibility and the solvability of matrix-valued asymmetric truncated Toeplitz operators. Specifically, in Theorem~\ref{G}, we provide a short algebraic proof of a characterization of 
 the invertibility of matrix-valued asymmetric truncated Toeplitz operators using Toeplitz
operators with $2 \times 2$ block matrix symbol. Finally, in Theorem~\ref{ATTO}, we discuss the normal solvability of matrix-valued asymmetric truncated Toeplitz operators. 

%The paper is organized as follows. In Section~\ref{s2}, after giving some introductory materials,  we give  an application of the matrix-valued truncated Toeplitz operators to prove the matrix version of Nevanlinna-Pick interpolation problem. In Section~\ref{s3}, we characterize the unitary and completely non-unitary parts of matrix-valued truncated Toeplitz operators. In Section~\ref {s4}, we provide a short algebraic proof of a characterization of  the invertibility of matrix-valued asymmetric truncated Toeplitz operators using Toeplitz operators with $2 \times 2$ block matrix symbol. Furthermore, we discuss the normal solvability of matrix-valued asymmetric truncated Toeplitz operators. % is also given.



%----------------------------------------------------
%----------------------------------------------------

%\ni  Let $\mathcal{L}_{n}$ denote the set of all $n\times n$ lower triangular block Toeplitz matrices.

%In analogy to the scalar-valued case, an element $\Theta\in H^2(\LL(E))$ is said to be an \emph{inner function} if its boundary values are almost everywhere unitary operators in $\LL(E)$.
%Let $$\mathcal{L}={\{\Theta\in H^{\infty}(\LL(E)): \Theta ~\text{is inner function}}\}$$
%and
%$$\mathcal{B}={\{B\in \mathcal{L}: B ~\text{is finite Blaschke-Potapov product}}\}.$$
%Note that 
%An inner function could be viewed as a power series as well as a Toeplitz operator.
%The space $\LL(E)^{n}$ is defined
%$$\LL(E)^{n}={\{(A_{0}, A_{1}, ..., A_{n-1}): A_{k}\in \LL(E)}\}.$$

%Let $W\in \LL(E)$ be a contraction. Then the operators defined as
%$$D_{W}=(I-W^{*}W)^{1/2} \quad\text{ and }\quad D_{W^{*}}=(I-WW^{*})^{1/2}$$ are called the \textit{defect operators}. The defect spaces of $W$ are $\mathcal{D}_{W}=\overline{D_{W}}$ and $\mathcal{D}_{W^{*}}=\overline{D_{W^{*}}}$. 
%For $\Theta\in \mathcal{L}$, the function $F_{W}(\Theta)(\lambda)\in \mathcal{L}(\mathcal{D}_{W}, \mathcal{D}_{W^{*}})$ defined by
%\begin{equation}\label{NTheta}
%F_{W}(\Theta)(\lambda)=\widetilde{\Theta}(\lambda)=-W+ D_{W^{*}}(I-\Theta(\lambda) W^{*})^{-1}\Theta(\lambda) D_{W}
%\end{equation}
%is also an inner function (see \cite{RK}), and $F_{W}$ is zero in $\mathcal{D}_{W}^{\perp}$.

%For each $n\in\mathbb{N}$, define $P_{n}: H^{2}(\LL(E))\to \mathcal{L}(E)^{n}$ by
%$$P_{n}(f)=P_{n}\left(\sum\limits_{k=0}^{\infty}A_{k}z^{k}\right)=(A_{0}, A_{1},...,A_{n-1})\quad \text{for every}\quad f\in H^{2}(\LL(E)).$$
%Moreover, define the operator $Q_{n}: H^{2}(\LL(E))\to \mathcal{L}_{n}$ by
%$$Q_{n}(f)=Q_{n}\left(\sum\limits_{k=0}^{\infty}A_{k}z^{k}\right)=\begin{bmatrix}
%A_{0} & 0 &0 &\cdots & 0 \\
%A_{1} & A_{0} & 0&\cdots & 0 \\
%A_{2} & A_{1} &A_{0}& \cdots & 0 \\
%\vdots & \vdots & \vdots &\ddots & \vdots \\
%A_{n-1} & A_{n-2} &A_{n-3} &\cdots & A_{0}
%\end{bmatrix}.$$
%There is a bijection $H:\LL(E)^{n}\to\mathcal{L}_{n}$ defined by
%$$H(A_{0}, A_{1},...,A_{n-1})=\begin{bmatrix}
%A_{0} & 0 &0 &\cdots & 0 \\
%A_{1} & A_{0} & 0&\cdots & 0 \\
%A_{2} & A_{1} &A_{0}& \cdots & 0 \\
%\vdots & \vdots & \vdots &\ddots & \vdots \\
%A_{n-1} & A_{n-2} &A_{n-3} &\cdots & A_{0}
%\end{bmatrix}$$
%such that
%$$Q_{n}=HP_{n}.$$
%\medskip
%\section {\bf{Blaschke-Potapov product}}\label{s3}
%Let $\Theta\in \mathcal{L}$ be a finite Blaschke-Potapov product as in \eqref{BP}. Then we have the following proposition.
%\begin{prop}
%The function $\widetilde{\Theta}$, defined in terms of an inner function $\Theta$ (as in \eqref{NTheta}) and a contraction $W$, given by
%\begin{equation}
%F_{W}(\Theta)(\lambda)=\widetilde{\Theta}(\lambda)=-W+ D_{W^{*}}(I-\Theta(\lambda) W^{*})^{-1}\Theta(\lambda) D_{W},
%\end{equation}
%is a finite Blaschke-Potapov product of degree $n$ if and only if $\Theta$ is a finite Blaschke-Potapov product of degree $n$.
%\end{prop}
%\begin{proof}
%Let $\Theta$ is a finite Blaschke-Potapov product of degree $n$. Then $\Theta$ is a rational inner function. Since the transformation
%$F_{W}$ preserves inner functions and rational functions (along with their degrees),  $F_{W}(\Theta)$  is a rational inner function of the same degree.  Furthermore, every rational inner   function can be
%represented as a Blaschke-Potapov product. Therefore, $\widetilde{\Theta}$ is also a finite Blashcke-Potapov product of degree $n$. The converse is proved similarly.
%\end{proof}
%Denote by $$\mathcal{L}_{n}^{c}={\{T\in \mathcal{L}_{n}: \|T\|\leq 1}\}.$$ The operator
%$\widetilde{F}_{W}:\mathcal{L}_{n}^{c}\to \mathcal{L}_{n}^{c}$ defined by
%$$\widetilde{F}_{W}(T)=-W+ D_{W^{*}}(I-%TW^{*})^{-1}TD_{W},$$ is well defined (see \cite{DT}).
%\begin{lem} \cite[Lemma 3]{DT}
%Let $T\in \mathcal{L}_{n}^{c}$. Then $\|\widetilde{F}_{W}(T)\|=1$  if and only if $\|T\|=1$ and $\|\widetilde{F}_{W}(T)\|<1$
 % if and only if $\|T\|<1$.
%\end{lem}

%\begin{rem}\cite[Lemma 3]{DT}
%The inverse of $\widetilde{F}_{W}$ is given by the formula
%\[
%(\widetilde{F}_{W}(T^{'}))^{-1}=W+ D_{W^{*}}(I+T^{'} W^{*})^{-1}T^{'} D_{W}.
%\]
%for any $T'\in \mathcal{L}_{n}^{c}$.
%\end{rem}
%Moreover, $F_{W}$ is a bijection since $$F_{W}^{-1}(\widetilde{\Theta})=F_{-W}(\widetilde{\Theta})=\Theta,$$ and
%$$\widetilde{F}_{W}^{-1}=\widetilde{F}_{-W}.$$
%The following diagram is commutative
%\[
%\begin{tikzcd}
%\mathcal{D}_{W} \arrow[r, "F_{W}"]\arrow[d, "Q_{n}"]& \mathcal{D}_{W^{*}}\arrow[d, "Q_n"] \\
%\mathcal{L}_{n}^{c}\arrow[r, "\widetilde{F}_{W}" ]& \mathcal{L}_{n}^{c}.
%\end{tikzcd}
%\]
%The following theorem is true for finite Blashcke product.
%\begin{thm}\cite [Theorem 2.3]{YAO}\label{unit norm}
%Let $b$ be a finite Blaschke product of degree $m-1$. Then for any $n\geq m$,
%$$\|Q_{n}(b)\|=1.$$
%\end{thm}

%We now give an example showing that Theorem~\ref{unit norm} may fail to be valid %is not true
% for a finite Blaschke-Potapov product B.

%\begin{ex}
%Consider
%$$B(z)=\begin{bmatrix}
%z & 0  \\
%0 & z
%\end{bmatrix}=zP+(I-P),
%$$
%where $P=\begin{bmatrix}
%1&0\\
%0&1
%\end{bmatrix}.$ Since $\text{degree}\,B=\text{rank}\,P=2,$ $B$ is a Blaschke-Potapov product of degree $2$. %, because $$\text{degree}~ B_{1}=\text{rank}~ P_{1}=2.$$
%Furthermore,
%$$B(z)=Pz=0+Pz+0\cdot z^{2}+0\cdot z^{3}+\cdots \,.$$
 %Now let $m-1=2$. Then $m=3$. Take $n=2<m=3$.
%Then
%$$Q_{2}(B)=\begin{bmatrix}
%\bf0 & \bf0  \\
%P & \bf0
%\end{bmatrix}=\begin{bmatrix}
%0 & 0&0&0  \\
%0&0&0&0\\
%1&0&0&0\\
%0&1&0 &0
%\end{bmatrix}.$$
%Hence, $\|Q_{2}(B)\|=1$.
%\end{ex}


%-------------------------------------------------------------------------------------------------------------


\section{\bf{Applications of matrix-valued} truncated Toeplitz operators}\label{s2}
In this section, we give an application of the matrix-valued truncated Toeplitz operators to prove the matrix version of Nevanlinna-Pick interpolation problem. 

The following lemma characterizes the finite dimensional vectorial model spaces.

 \begin{lem}\cite[Lemma 5.1]{Peller}
Let $\mathcal{M}$ be an invariant subspace of the operator $S$ (namely, multiplication by $z$ on $H^2(E)$). Then $\mathcal{M}$ has finite codimension if and only if $\mathcal{M}=\Theta H^2(E)$ for a Blaschke-Potapov product $\Theta$ of finite degree. Moreover,
$$\text{dim}K_{\Theta}=\text{codim}\mathcal{M}=\text{deg}~ \Theta.$$
 \end{lem}

For each $\lambda\in \mathbb{D}$ and $x\in E$, the function
$$k_{\lambda}^{\Theta}x=\frac{1}{1-\bar \lambda z}(I_{E}-\Theta(z)\Theta(\lambda)^*)x$$
belong to $K^{\infty}_{\Theta}=K_{\Theta}\cap H^{\infty}(E)$ and has the
 reproducing property
$$\langle f, k_{\lambda}^{\Theta}x\rangle=\langle f(\lambda), x\rangle \quad \text{for every}\quad f\in K_\Theta.$$
When $\Theta$ is a pure inner function, i.e., $\|\Theta(0)\|< 1$, 
$$k_{0}^{\Theta}(z)= I_{E}-\Theta(z)\Theta(0)^*,$$ which is invertible in $H^{\infty}(\oL)$. 


In the case when $\Theta(\lambda)=0$, we have
$$k_{\lambda}^{\Theta}x=k_{\lambda}x.$$
The \emph{model operator} $S_{\Theta}\in\mathcal{L}(K_{\Theta})$ is defined  by the formula
\begin{equation*}\label{eq:Model Operator}
(S_{\Theta}f)(z)=P_{\Theta} (zf), \quad f\in K_{\Theta}.
 \end{equation*}
Its adjoint is given by
\[
(S_{\Theta}^{*}f)(z)=\frac{f(z)-f(0)}{z},
\]
which is the restriction of the left shift in $H^{2}(E)$ to the $S^{*}$-invariant subspace $K_{\Theta}$.

The space $\Theta H^2(E)$ is invariant under $S=M_z$, but it is not invariant under $M_\Phi$ for general $\Phi\in H^{\infty}(\oL)$, and hence $K_\Theta$ is not invariant under $M^*_\Phi$ (see \cite[Example 4.1]{RT}). Therefore, the relation $(A^{\Theta}_\Phi)^*=M^*_\Phi|_{K_\Theta}$ does not hold. However, if $\Phi\in H^{\infty}(\oL)$ commutes with $\Theta$, then $M^*_\Phi K_\Theta\subset K_\Theta$ and hence $(A^{\Theta}_\Phi)^{*}=M^*_\Phi|_{K_\Theta}$. It then follows that $(A^{\Theta}_\Phi)^{*} S^*_\Theta=S^*_\Theta (A^{\Theta}_\Phi)^*$, which gives $A^{\Theta}_\Phi S_\Theta=S_{\Theta} A^{\Theta}_\Phi$.

\medskip

The following theorem shows that the converse is also true.

%{\color{red} Actually, Theorem 2.2 is the commutant lefting theorem.} 

\begin{thm}\cite[Chapter VI]{NF}\label{clt}
    Let $\Theta$ be an inner function, and let $A\in \mathcal{L}(K_{\Theta})$ such that $AS_{\Theta}=S_{\Theta}A$. Then there exists $\Phi\in H^\infty(\mathcal{L}(E))$ with $\Phi\Theta=\Theta\Phi$ such that $A=A_{\Phi}^{\Theta}$. Moreover %we have:
    \begin{enumerate}
        \item For any $A=A_{\Phi}^{\Theta}$, with $\Phi\in H^{\infty}(\oL)$, we have
        $$\|A\|=dist(\Phi, \Theta H^\infty(\oL)$$ as well as %and
        $$\|A\|=\inf{\{\|\Phi\|_{\infty}: A=A_{\Phi}^{\Theta},~~\Phi\in H^{\infty}(\mathcal{L}(E))}\}.$$
        
        \item There exists $\Phi\in H^\infty(\oL)$ such that $\|A\|=\|\Phi\|_{\infty}.$
    \end{enumerate}
\end{thm}
\medskip

To prove the matrix version of the Nevanlinna-Pick interpolation problem, we recall some needed notions.

Given square matrix $A=(a_{ij})$, by the notation $A\geq 0$ we mean that
$$\sum_{i,j=1}^{n}\overline{\alpha}_i\alpha_j a_{i,j}\geq 0$$
 for every $\alpha_1, \alpha_2,\dots,\alpha_n \in \mathbb{C}$. 
 If so, we say that $A$ is positive definite. 

Let $E$ be a separable complex Hilbert space. A function $K: \mathbb{D} \times \mathbb{D}\to \oL$ is called a \emph{kernel function} provided that for every choice of $n$ distinct points $\lambda_1, \lambda_2,.\dots,\lambda_n$ in $\mathbb{D}$ and any vectors $x_1,x_2,\dots,x_n$ in $E$, we have
$$\sum_{i,j=1}^n\big\langle K(\lambda_i, \lambda_j)x_j, x_i\big\rangle\geq 0.$$

\begin{prop}\cite[Proposition 2.13]{VP}\label{VP}
     Let $E$ be a reproducing kernel Hilbert space with reproducing kernel $K$. Then $K$ is a kernel function.
\end{prop}

 We shall make use of the following lemma.

\begin{lem}
    Let $\Phi\in H^\infty(\oL)$ and $\lambda\in \mathbb{D}$ such that $\Theta(\lambda)=0$. Then 
    $$A_{\Phi^*}^{\Theta}k_{\lambda}^{\Theta}x=\Phi(\lambda)^*k^{\Theta}_{\lambda}x.$$
\end{lem}

\begin{proof}
Let $f\in K_{\Theta}$ and $x\in E$. Then
\begin{align*}
    \langle A_{\Phi^*}^{\Theta}k_{\lambda}^{\Theta}x, f\rangle
   & =\langle \Phi^* k_{\lambda}^{\Theta}x, f\rangle\\
    &=\langle  k_{\lambda}^{\Theta}x, \Phi f\rangle\\
      &=\overline{\langle  \Phi f, k_{\lambda}^{\Theta}x\rangle}\\
      &=\overline{\langle  \Phi(\lambda) f(\lambda), x\rangle}\\
      &=\overline{\langle   f(\lambda), \Phi(\lambda)^* x\rangle}\\
&=\overline{\langle   f, k_{\lambda}^{\Theta}\Phi(\lambda)^* x\rangle}\\
&=\overline{\langle   f,\Phi(\lambda)^* k_{\lambda}^{\Theta} x\rangle}\\
&=\langle \Phi(\lambda)^* k_{\lambda}^{\Theta} x,  f\rangle.
\end{align*}
The required result follows.
\end{proof}

As an application of matrix-valued truncated Toeplitz operators, we discuss a matrix version of the Nevanlinna–Pick interpolation problem, which  addresses the following question.

    Let $(\lambda_i)^{d}_{i=1}\subset \mathbb{D}$ and $(W_i)_{i=1}^d\subset \oL$ with $W_j^{*}W_i=W_iW_j^{*}$. Does there exist a function $\Phi\in H^\infty(\oL)$ such that
    \begin{equation}
        \Phi(\lambda_i)=W_i, \quad \text{for}~~ 1\leq i\leq d, \quad \text{and} \quad \|\Phi\|_\infty\leq 1?\nonumber
    \end{equation}
%\end{thm}
If so, what is the characterizing condition? 

The answer is positive and the characterizing condition for the problem to have a solution is similar to that of the classical version of the theorem. 

\begin{thm}\label{NevPick}
     Let $(\lambda_i)^{d}_{i=1}\subset \mathbb{D}$ and $(W_i)_{i=1}^d\subset \oL$ such that $W_j^{*}W_i=W_iW_j^{*}$ for $1\leq i,j\leq d$. Then the following are equivalent:
     \begin{enumerate}
         \item There exist a function $\Phi\in H^\infty(\oL)$ such that \begin{equation*}
        \Phi(\lambda_i)=W_i, \quad \text{for}~~ 1\leq i\leq d, \quad \text{and} \quad \|\Phi\|_\infty\leq 1.\label{i}
    \end{equation*}
    \item The block matrix
    $$P=\Big[\frac{I-W^*_iW_j}{1-\overline{\lambda_i}\lambda_j}\Big]_{i,j=1}^{d}$$ is positive definite. \label{ii}
    \end{enumerate}
\end{thm}

\begin{proof}
Let $\Theta$ be a finite \emph{Blaschke-Potapov} product such that $\Theta(\lambda_i)=0$ for all $1\leq i\leq d$. Then the model space corresponding to $\Theta$ is given by
$$K_\Theta=\text{span}{\{k_{\lambda_i}x: \lambda_i\in \mathbb{D}}\}.$$
Define $R: K_{\Theta}\to K_{\Theta}$ by
$$Rk_{\lambda_i}x=W_{i}^*k_{\lambda_i}x.$$
Then
\begin{equation}\label{com1}
RS^*_{\Theta}k_{\lambda_i}x=R\overline{\lambda}_ik_{\lambda_i}x=W_{i}^*\overline{\lambda}_ik_{\lambda_i}x
\end{equation}
and
\begin{equation}\label{com2}
S^*_{\Theta}Rk_{\lambda_i}x=S^*_{\Theta}W^*_ik_{\lambda_i}x=\overline{\lambda}_iW^*_ik_{\lambda_i}x=W_{i}^*\overline{\lambda}_ik_{\lambda_i}x.
\end{equation}
From \eqref{com1} and \eqref{com2} we have $RS^*_{\Theta}=S^*_{\Theta}R$. By Theorem \ref{clt}, there exists $\Phi\in H^\infty(\oL)$ with $\Phi\Theta=\Theta\Phi$ such that $R=A_{\Phi^*}^{\Theta}$  and
$$Rk_{\lambda_i}x=A^\Theta_{\Phi^*}k_{\lambda_i}x=\Phi(\lambda_i)^*k_{\lambda_i}x.$$
Hence $\Phi(\lambda_i)=W_i$ for all $1\leq i\leq n$.
Furthermore, we also have
\begin{equation}\label{rphi}
\|R\|=\|\Phi\|_{\infty}.
\end{equation}
Let $x_1, x_2,...,x_d\in E$. Then
$$
X=\begin{bmatrix}
    x_1\\
    x_2\\
    \vdots\\
   % .\\
   % .\\
   % .\\
    x_d
\end{bmatrix}\in E\oplus E\oplus...\oplus E.
$$
Applying the block matrix $P$, we have
$$PX=\begin{bmatrix}
   \frac{I-W^*_1W_1}{1-\overline{\lambda_1}\lambda_1}x_1+ \frac{I-W^*_1W_2}{1-\overline{\lambda_1}\lambda_2}x_2+\dots+\frac{I-W^*_1W_d}{1-\overline{\lambda_1}\lambda_d}x_d \\
    \frac{I-W^*_2W_1}{1-\overline{\lambda_2}\lambda_1}x_1+ \frac{I-W^*_2W_2}{1-\overline{\lambda_2}\lambda_2}x_2+\dots+\frac{I-W^*_2W_d}{1-\overline{\lambda_2}\lambda_d}x_d\\
    \vdots\\
    \frac{I-W^*_dW_1}{1-\overline{\lambda_d}\lambda_1}x_1+ \frac{I-W^*_dW_2}{1-\overline{\lambda_d}\lambda_2}x_2+\dots+\frac{I-W^*_dW_d}{1-\overline{\lambda_d}\lambda_d}x_d
    \end{bmatrix}.$$
For $i,j=1,\dots,d$, using the assumption $W_j^*W_i=W_i^*W_j$,  we have
\ben \langle W_j^*W_i k_{\lambda_j}x_i, k_{\lambda_i} x_j\rangle&=&\langle W_i^*W_jk_{\lambda_j}x_i, k_{\lambda_i} x_j\rangle\nn\\
&=&\langle W_i^*R k_{\lambda_j} x_i,k_{\lambda_i} x_j\rangle\nn\\
&=&\langle R k_{\lambda_j}x_i, W_i k_{\lambda_i} x_j\rangle\nn\\
&=&\langle R k_{\lambda_j}x_i, R k_{\lambda_i} x_j\rangle\nn\\
&=&\langle R^*R k_{\lambda_j}x_i, k_{\lambda_i} x_j\rangle.\nn\eeqn  
We obtain the following expression of the inner product $\langle PX, X\rangle$ in terms of the  matrix with entries 
$k_{\lambda_j(\lambda_i)}$, \, $i,j=1,\dots,d$.
 
\begin{align*}
    \langle PX, X\rangle 
%&=\Biggl\langle\begin{bmatrix}
 %  \frac{I-W^*_1W_1}{1-\overline{\lambda_1}\lambda_1}x_1+ \frac{I-W^*_1W_2}{1-\overline{\lambda_1}\lambda_2}x_2+...+\frac{I-W^*_1W_d}{1-\overline{\lambda_1}\lambda_d}x_d \\
 %   \frac{I-W^*_2W_1}{1-\overline{\lambda_2}\lambda_1}x_1+ \frac{I-W^*_2W_2}{1-\overline{\lambda_2}\lambda_2}x_2+...+\frac{I-W^*_2W_d}{1-\overline{\lambda_2}\lambda_d}x_d\\
 %   \vdots\\
%    \frac{I-W^*_dW_1}{1-\overline{\lambda_d}\lambda_1}x_1+ \frac{I-W^*_dW_2}{1-\overline{\lambda_d}\lambda_2}x_2+...+\frac{I-W^*_dW_d}{1-\overline{\lambda_d}\lambda_d}x_d
%    \end{bmatrix}
%,\begin{bmatrix}
%    x_1\\
%    x_2\\
  %  .\\
%    .\\
  %  .\\
  %  x_d
%\end{bmatrix}\Biggr\rangle%\displaybreak
%\\
%
&=\sum_{i,j=1}^d\Big\langle \frac{I-W^*_jW_i}{1-\overline{\lambda}_j\lambda_i}x_i, x_j\Big\rangle\\
&=\sum_{i,j=1}^d\langle (I-W^*_jW_i)k_{\lambda_j}(\lambda_i)x_i, x_j\rangle\\
&=\sum_{i,j=1}^d\langle (I-W^*_jW_i)k_{\lambda_j}x_i, k_{\lambda_i}x_j\rangle\\
&=\sum_{i,j=1}^d\langle (I-R^*R)k_{\lambda_j}x_i, k_{\lambda_i}x_j\rangle\\
&=\sum_{i,j=1}^d\langle (I-R^*R)k_{\lambda_j}(\lambda_i)x_i, x_j\rangle.
\end{align*}
Since $k_{\lambda}x$ is a reproducing kernel function for the model space $K_\Theta$, by Proposition~\ref{VP} it follows that the matrix
$$[k_{\lambda_j}(\lambda_i)]_{i,j=1}^d$$
is positive definite.
\medskip

%$1. \Longleftrightarrow 2.$
Noting that 
\begin{equation}\label{R}
P\geq 0 \quad \iff \quad 1-R^*R\geq 0 \quad  \iff\quad \|R\|\leq 1,
\end{equation}
from \eqref{rphi} and \eqref{R} we have $\|\Phi\|_{\infty}\leq 1$, which shows the equivalence of statements (1) and (2). 
\end{proof}

 \section{\bf{Unitary and completely non-unitary parts of MTTO's}}\label{s3}

Recall that, given a complex Hilbert space $E$,  an operator $T\in \oL$ is said to be a contraction if $\|T\|\leq 1$. An operator $T\in \oL$ is said to be \emph{completely non-unitary} (c.n.u) if there does not exist a $T$-reducing subspace $\mathcal{M}\subset E$ such that $T|_{\mathcal{M}}$ is unitary. The following result in \cite{NF} gives a decomposition of a contraction $T\in \oL$ into a direct sum of a unitary and a completely non-unitary operators on certain closed $T$-reducing subspaces of $E$.

\medskip


 \begin{thm}\cite{NF}\label{ucnu}
     Let $T\in \oL$ be a contraction.  Then $E=\h_u\oplus \h_{c.n.u}$ where $\h_u,\h_{c.n.u}\subset E$ are closed $T$-reducing subspaces of $E$ and $T$ can be decomposed as
     $$T=T|_{\hsp_u}\oplus T|_{\hsp_{c.n.u}}$$ where $T|_{\h_u}$ is unitary and $T|_{\h_{c.n.u}}$ is completely non-unitary.
 \end{thm}

\medskip

 
 Recall that an inner function $\Theta$ with $\|\Theta(0)\|<1$ is called pure. Throughout
this section we assume $\Theta$ to be a pure inner function.

\medskip



%{\color{red} Reading your paper with Timotin in CAOT, I noticed that you assumed $\Theta$ to be purely inner before mentioning $k_\lambda^\Theta(z)$ (see p. 1000). Why not in this paper? In Theorem 3.4, if $\Theta_j$ ($j=1,2$) are pure inner functions and $\Theta=\Theta_1\Theta_2$, then $\Theta$ is also pure. Is the converse true? The assumptions  in the statement of the theorem should be clarified before making a blanket statement above "all inner functions in this section are assumed to be pure".} 

%{\color{orange} The converse is not true. Assume 
%$$\Theta_1(z)=\begin{pmatrix}
%   1&0\\
%    0& \frac{z-\frac{1}{2}}{1-\frac{1}{2}z}
%\end{pmatrix}. $$ Then $\Theta_1$ is an inner function, and
%Similarly.
%$$\Theta_2(z)=\begin{pmatrix}
%  \frac{z-\frac{1}{3}}{1-\frac{1}{3}z}&0\\
 %   0& 1
%\end{pmatrix}. $$ Then $\Theta_2$ is also an inner function, and  
%$$\|\Theta_{2}(0)\|=\text{max}\left\{1, \frac{1}{3}\right\}=1.$$
%But
%$$\Theta(z)=\Theta_1\Theta_2=\begin{pmatrix}
 %  \frac{z-\frac{1}{3}}{1-\frac{1}{3}z}&0\\
 %   0& \frac{z-\frac{1}{2}}{1-\frac{1}{2}z}
%\end{pmatrix}. $$ Then 
%$$\|\Theta(0)\|=\text{max}\left\{\frac{1}{2}, \frac{1}{3}\right\}=\frac{1}{2}.$$
%So $\Theta$ is pure.}


Let $\Phi \in L^{\infty}(\mathcal{L}(E))$ such that $\|\Phi\|_{\infty} \leq 1$. Then, $A^{\Theta}_{\Phi}$ is a contraction, since by part (1) of Theorem~\ref{clt}, 
\[
\|A^{\Theta}_{\Phi} \| \leq \|\Phi\|_{\infty} \leq 1.
\]
The unitary subspace of $A^{\Theta}_{\Phi}$ is defined as
\[
\mathcal{H} = \{f \in K_\Theta: \|(A^{\Theta}_{\Phi})^n f\| = \|f\| = \|(A^{\Theta}_{\Phi})^{*n} f\| \text{ for all }\, n \in \mathbb{N} \}.
\]
In particular, for $f \in \h$, we have
\[
\|f\| = \|A^{\Theta}_{\Phi}f\| = \|P_{\Theta} (\Phi f) \| \leq \|\Phi f \| \leq \|f\|.
\]
Hence
\[
\|P_{\Theta} (\Phi f) \| = \|\Phi f \|,
\]
which implies that $\Phi f \in K_\Theta$.  Moreover, $\|\Phi f\| = \|f \|$. In a similar fashion, using the condition $\|(A^{\Theta}_{\Phi})^{*}f \| = \|f\|$,  we obtain  $\Phi^* f \in K_\Theta$ and $\|\Phi^* f\| = \|f \|$. The same property holds for all powers $n \in \mathbb{N}$. Hence,
\[
\mathcal{H} = \{f \in K_\Theta: \Phi^{n} f, \Phi^{*n} f \in K_\Theta \text{ and } \|\Phi^n f\| = \|f\| = \|\Phi^{*n} f\| \text{ for all }\, n \in \mathbb{N} \}.
\]


\noindent Using the multiplication operator notation,
\[
\mathcal{H} = \{f \in K_\Theta: L_{\Phi}^n f, L_{\Phi}^{*n} f \in K_\Theta \text{ and } \|L_{\Phi}^n f\| = \|f\| = \|L_{\Phi}^{*n} f\| \text{ for all }\, n \in \mathbb{N} \}.
\]
 Furthermore, since $L_{\Phi}$ is a contraction, for any $n \in \mathbb{N}$, we have
 \[
\|L_{\Phi}^n f \| = \|f\| \Longleftrightarrow f \in  \ker (I - L_{\Phi}^{*n} L_{\Phi}^{n}),
\]
and
\[
\|L_{\Phi}^{*n} f \| = \|f\| \Longleftrightarrow f \in  \ker (I - L_{\Phi}^{n} L_{\Phi}^{*n}).
\]
Thus, we may express $\h$ as
\begin{equation}
\mathcal{H} = \{f \in K_\Theta: L_{\Phi}^n f, L_{\Phi}^{*n} f \in K_\Theta \text{ and } L_{\Phi}^n L_{\Phi}^{*n} f = f = L_{\Phi}^{*n} L_{\Phi}^n f \text{ for all }\, n \in \mathbb{N} \}.\label{calHform}\end{equation}
We deduce the following Lemma.
\medskip


\begin{lem}\label{unitary part}
     Let $f\in K_\Theta$ such that $L^n_{\Phi}f, L^{*n}_{\Phi}f\in K_\Theta$. Then $f\in\h$  if and only if $\Phi(\zeta)$ is unitary $\text{a.e. on } ~~\mathbb{T}.$
\end{lem}

\begin{proof} Let $f\in \hsp$. Then from the equality
$$L_{\Phi}^n L_{\Phi}^{*n} f = f = L_{\Phi}^{*n} L_{\Phi}^n f,$$ it follows that
$$\Phi(\zeta)\Phi^{*}(\zeta)f(\zeta)=f(\zeta)=f(\zeta)\Phi^{*}(\zeta)\Phi(\zeta).$$
\vskip3pt

\noindent Thus, for almost every $\zeta\in \mathbb{T}$, %we have
 either $f(\zeta)=0$ or $\Phi(\zeta)$ is unitary.  Since $f$ is analytic, it cannot be zero constant on a set of positive measure. Therefore, $\Phi$ must be unitary a.e. on $\mathbb{T}$. The converse is straightforward.
\end{proof}



\begin{cor}
The model space $K_{\Theta}$ has a unitary part if and only if $K_{\Theta}$ is reducing under $L_{\Phi}$ and there exists a subset $K$ of $\mathbb{T}$ of positive measure such that $\Phi(\zeta)$ is unitary on $K$.
\end{cor}
%\begin{proof}
The proof follows from Lemma~\ref{unitary part}.
%\end{proof}

\begin{thm}\label{Stheta}
     Every $S_\Theta$-invariant subspace of $K_\Theta$ is of the form
     $\Theta_{1} K_{\Theta_2}$ where $\Theta_1$ and $\Theta_2$ are inner functions such that $\Theta=\Theta_1\Theta_2$.
\end{thm}

\begin{proof}
Let $\mathcal{M}\subset K_{\Theta}$ be a closed subspace such that $S_{\Theta}\mathcal{M}\subset \mathcal{M}$. Define $\mathcal{N}$ as the direct sum
$$\mathcal{N}=\mathcal{M}\oplus \Theta H^{2}(E).$$ We now show that $\mathcal{N}$ is invariant under the forward shift $S$. Take $f\in \mathcal{N}$ and let $g\in \mathcal{M}$ and $h\in H^{2}(E)$ such that
$$f=g+\Theta h\,.$$ %for some $g\in \mathcal{M}$ and $h\in H^{2}(E)$. 
Applying the forward shift $S$, we have 
\begin{eqnarray} Sf=zf=zg+z\Theta h=Sg+\Theta zh,\label{Sf1}\end{eqnarray}
where we note that the second term on the right-hand side is in $\Theta H^{2}(E)$. For the first term, we note that since $g\in \mathcal{M}\subset K_{\Theta}$, we have
%$S_{\Theta}g\in \mathcal{M}$ and 
$$S_{\Theta}g=P_{\Theta}(zg)=(I-P_{\Theta H^{2}(E)})(zg)=Sg-P_{\Theta H^{2}(E)}(zg),$$ which gives
$$Sg=S_{\Theta}g+P_{\Theta H^{2}(E)}(zg).$$
Since $\mathcal{M}$ is $S_{\Theta}$ invariant,  $g'=S_{\Theta}g$ is in $\mathcal{M}$ and 
\begin{eqnarray} Sg=g'+\Theta h'\label{Sg2}\end{eqnarray} for some %$g'\in \mathcal{M}$ and
 $h'\in H^{2}(E).$
Substituting (\ref{Sg2}) into (\ref{Sf1}) yields
$$Sf=%g'+\Theta h'+\Theta zh=
g'+\Theta(h'+zh)%=g'+\Theta k
\in \mathcal{M}\oplus\Theta H^{2}(E)=\mathcal{N}\,.$$
%where $k=h'+zh\in H^{2}(E)$. 
Therefore $\mathcal{N}$ is $S$-invariant. By  Beurling-Lax-Halmos theorem (see \cite{AM}), there exists an inner function $\Theta_{1}$ such that $$\mathcal{N}=\Theta_{1}H^{2}(E).$$
We also have $\Theta H^{2}(E)\subset \mathcal{N}=\Theta_{1}H^{2}(E)$. Thus there exist an analytic function $\Theta_{2}\in H^{\infty}(\mathcal{L}(E))$ such that $$\Theta=\Theta_{1}\Theta_{2}.$$ 
The function $\Theta_2$ is analytic (since $\Theta_{1}$ is a factor of $\Theta$) and can be expressed as
%We can write $\Theta_2$ as 
$$\Theta_{2}=\Theta_1^{*}\Theta.$$ Then
$$\Theta_{2}^{*}\Theta_2=(\Theta_1^{*}\Theta)^{*}(\Theta_1^{*}\Theta)=\Theta^*\Theta_1\Theta_{1}^{*}\Theta=I.$$ 
Thus $\Theta_{2}$ is an inner function. Writing $\mathcal{M}$ as
$N\ominus \Theta H^{2}(E),$ it follows that
$$\mathcal{M}=\Theta_{1}H^{2}(E)\ominus \Theta_{1}\Theta_{2}H^{2}(E)=\Theta_{1}[H^{2}(E)\ominus \Theta_{2}H^{2}(E)]=\Theta_{1}K_{\Theta_2},$$
as desired.
\end{proof}
%\smallskip

\begin{thm}\label{inv}
    The unitary space $\h$ is $S_\Theta$-invariant if $L_{\Phi}K_{\Theta}\subseteq K_{\Theta}$.
\end{thm}

\begin{proof}  The assumption $L_{\Phi} K_{\Theta}\subseteq K_\Theta$ is equivalent to $$L_{\Phi}P_{\Theta}=P_{\Theta}L_{\Phi}.$$
    Let $f\in \h$. Then, by (\ref{calHform}), we obtain %    $L_{\Phi}^n f, L_{\Phi}^{*n} f\in K_\Theta$ and $ L_{\Phi}^n L_{\Phi}^{*n} f = f = L_{\Phi}^{*n} L_{\Phi}^n f,$ for all $n \in \mathbb{N}.$    Observe that
$$L_\Phi S_\Theta f=L_\Phi P_\Theta (zf)=P_\Theta L_\Phi L_z f=P_\Theta L_z L_\Phi f=S_{\Theta}L_\Phi f\in K_\Theta.$$
Repeating this process, we see that $S_\Theta L^n_\Phi=L^{n}_{\Phi}S_\Theta$ for all $n\in \mathbb{N}$. Similarly, we have $S_\Theta L^{*n}_\Phi=L^{*n}_{\Phi}S_\Theta$.
Hence,
    $$L^n_{\Phi}S_\Theta f=S_\Theta L^n_\Phi f\in K_\Theta$$ and
    $$L^{*n}_{\Phi}S_\Theta f=S_\Theta L^{*n}_\Phi f\in K_\Theta.$$
     Therefore, we have
    $L_{\Phi}^n S_\Theta f, L_{\Phi}^{*n} S_\Theta f\in K_\Theta$ for all $f\in \h$ and for all $n\in \mathbb{N}$.

    Furthermore, we have
    $$L^n_{\Phi}L^{*n}_{\Phi}S_\Theta f=S_\Theta L^n_{\Phi}L^{*n}_{\Phi} f=S_\Theta f$$
    and
    $$L^{*n}_{\Phi}L^{n}_{\Phi}S_\Theta f=S_\Theta L^{*n}_{\Phi}L^{n}_{\Phi} f=S_\Theta f.$$
This shows that $S_\Theta f\in \h$, hence $\h$ is $S_\Theta$-invariant.\end{proof}

We now provide a counterexample to show that the converse of Theorem \ref{inv} does not hold.
%{\color{blue}
\begin{ex}
Let $E=\mathbb{C}^{2}$ and consider
$$
\Theta(z)=
\begin{pmatrix}
z&0\\
0&z
\end{pmatrix}
=zI_{E}.
$$
%\end{ex}
The model space corresponding to $\Theta$ is
$$K_{\Theta}=\left\{\begin{pmatrix}
    a\\
    b
\end{pmatrix}:  a, b\in \mathbb{C}\right\}.$$
Let
$$
\Phi(z)=
\begin{pmatrix}
1&0\\
0&z
\end{pmatrix}.
$$
Then $\Phi\in L^{\infty}(\oL)$ and $\|\Phi\|_{\infty}=1$. Observe that $A_{\Phi}^{\Theta}$ acts on $K_{\Theta}$ as a matrix. Indeed for $f=\begin{pmatrix}
    a\\
    b
\end{pmatrix}\in K_{\Theta},$
$$A_{\Phi}^{\Theta}f=P_{\Theta}(\Phi f)=P_{\Theta}
\begin{pmatrix}
a\\
bz
\end{pmatrix}
=\begin{pmatrix}
a\\
0
\end{pmatrix}=\begin{pmatrix}
1&0\\
0&0
\end{pmatrix}f.$$
%Hence, $A_{\Phi}^{\Theta}$ acts on $K_{\Theta}$ as a matrix:
%$$A_{\Phi}^{\Theta}=\begin{pmatrix}
%1&0\\
%0&0
%\end{pmatrix}.$$
 Thus, $A_{\Phi}^{\Theta}$ is a projection. Its unitary space is
$$\hsp=\text{ker}(I-(A_{\Phi}^{\Theta})^{*}A_{\Phi}^{\Theta})\cap \text{ker}(I-A_{\Phi}^{\Theta}(A_{\Phi}^{\Theta})^{*}).$$
Since  $A_{\Phi}^{\Theta}(A_{\Phi}^{\Theta})^{*}=A_{\Phi}^{\Theta}=(A_{\Phi}^{\Theta})^{*}A_{\Phi}^{\Theta}$, it follows that %=A_{\Phi}^{\Theta}.$ So
$$\hsp=\text{ker}(I-A_{\Phi}^{\Theta})=\text{Rang} A_{\Phi}^{\Theta}=\text{span}\left\{\begin{pmatrix}
    1\\
    0
\end{pmatrix}%e_1
\right\}.$$

%where $e_1=\begin{pmatrix}
 %   1\\
   % 0
%\end{pmatrix}$.
%Hence, $$\hsp=\left\{\begin{pmatrix}
%    a\\
%    0
%\end{pmatrix}
%: a\in \mathbb{C}\right\}.$$
Observe that the model operator $S_{\Theta}$ is trivial on $K_{\Theta}$. 
Indeed, for $f=\begin{pmatrix}
    a\\
    b
\end{pmatrix}\in K_{\Theta}$, %the model operator is given by
$$S_{\Theta}f=P_{\Theta}(zf)=P_{\Theta}\begin{pmatrix}
    az\\
    bz
\end{pmatrix}=0.$$ %Thus $S_{\Theta}=0$ on $K_{\Theta}$. 
Therefore,
%$$S_{\Theta}\hsp=\left\{0\right\}\subset \hsp.$$ So
$\hsp$ is trivially $S_{\Theta}$-invariant. However, the inclusion $L_{\Phi}K_{\Theta}\subset K_{\Theta}$ does not hold. To see this, consider $f=\begin{pmatrix}
    1\\
    1
\end{pmatrix}\in K_{\Theta}$. Then
$$L_{\Phi}f=\Phi f=\begin{pmatrix}
    1&0\\
    0&z
\end{pmatrix}\begin{pmatrix}
    1\\
    1
\end{pmatrix}=\begin{pmatrix}
    1\\
    z
\end{pmatrix}\notin K_{\Theta}.$$ We conclude that
$$L_{\Phi}K_{\Theta}\not\subset K_{\Theta}.$$
\end{ex}
%\end{proof}
\medskip

The model space corresponding to the inner function $\Theta=\Theta_1 \Theta_2$ can be decomposed as $$K_\Theta=K_{\Theta_1}\oplus \Theta_1 K_{\Theta_2}.$$ By Theorem \ref{Stheta}, $\Theta_1 K_{\Theta_2}$ is invariant under $S_{\Theta}$. Its orthogonal complement $K_{\Theta_1}=[\Theta_1 K_{\Theta_2}]^{\perp}$ is invariant under  $S^*_\Theta$. 

    \medskip
By \cite{RT}, every matrix-valued truncated Toeplitz operator has a symbol in a certain class. Denote by $M_{\Theta}$ the orthogonal complement of $\Theta H^2(\oL)$ in $H^2(\oL)$ endowed with the Hilbert–Schmidt norm. %{\color{red} This sentence was taken from your paper with Timotin (top of p. 1001). Could you be more specific what is from [11] you are using?} 
%
%{\color{orange} I am using the analyticity if $\Phi$ and $\Psi$, and $A_{\Psi}^{\Theta}=A_{\Phi}^{\Theta}$, then reaching the conclusion that $\Psi=\Phi+k_{0}^{\Theta}X$.
%
%For example, if you look at Corollary 6.4 of my paper with Timotin, there are no $\Psi_2^{*}$ and $\Psi_2^{'*}$ (here in our case) because our symbols are bounded analytic here. If we consider $\Psi_2^{*}=\Psi_2^{'*}=0$ (because both of them are co-analytic), then we have $A=A_{\Psi_1}$ and $A=A_{\Psi_1^{'}}$. Hence, $A_{\Psi_1}=A_{\Psi_1^{'}}$ which follows that $\Psi_{1}=\Psi_{1}^{'}+k_{0}^{\Theta}X$.}
%
Moreover, $$L^2(\oL)=M_{\Theta}+\Theta H^2(\oL)+M_{\Theta}^{*}+[\Theta H^2(\oL)]^{*}.$$

\begin{thm}\label{Phi constant}
The model space $K_{\Theta}$ is invariant with respect to $L_{\Phi}$ if and only if $\Phi(\zeta)$ is constant.
\end{thm}

\begin{proof}
Assume $K_{\Theta}$ is invariant under $L_{\Phi}$. Then
$$A_{\Phi}^{\Theta}f=P_{\Theta}(\Phi f)=\Phi f=L_{\Phi}f.$$
Thus, $L_{\Phi}=A_{\Phi}^{\Theta}$.

Now we show that $\Phi$ is bounded and analytic. Take $f=k_{0}^{\Theta}x\in K_{\Theta}$ for any $x\in E$, we have 
$$\Phi(z)k_{0}^{\Theta}x=\Phi (I-\Theta(z)\Theta(0)^{*})x\in H^2(E).$$ Since $\Theta$ is pure, i.e., $\|\Theta(0)\|<1$, we have
$$(I-\Theta(z)\Theta(0)^{*})^{-1}\in H^{\infty}(\oL).$$ Therefore
$$\Phi x=\Phi (I-\Theta(z)\Theta(0)^{*})(I-\Theta(z)\Theta(0)^{*})^{-1}x\in H^2(\oL).$$ It follows that 
$$\Phi\in H^2(\oL)\cap L^{\infty}(\oL)=H^{\infty}(\oL).$$

Since $L_{\Phi}$ commutes with multiplication operator $L_{z}$, we have 
$$S_{\Theta}L_{\Phi}=L_{\Phi}S_{\Theta}.$$
By Theorem \ref{clt}, there exists $\Psi\in H^{\infty}(\mathcal{L}(E))$ such that $\Psi\Theta=\Theta \Psi$ and 
$$L_{\Phi}=A_{\Psi}^{\Theta}.$$
Hence,
$$A_{\Psi}^{\Theta}S_{\Theta}=S_{\Theta}A_{\Psi}^{\Theta}.$$
But since $L_{\Phi}=A_{\Phi}^{\Theta}$ on $K_{\Theta}$, we have 
\begin{equation}\label{MTTOEQ}
A_{\Phi}^{\Theta}=A_{\Psi}^{\Theta} 
\end{equation}
on $K_{\Theta}$.
Using the reproducing kernel property and invariance of $K_{\Theta}$ under $L_{\Phi}$, for each $\lambda\in \mathbb{D}$, $x\in E$, and $f\in K_\Theta$, we have 
\begin{align}
\langle (A_{\Phi}^{\Theta})^{*}k_{\lambda}^{\Theta}x, f\rangle
&=\langle k_{\lambda}^{\Theta}x, \Phi f\rangle\nn\\
&=\overline{\langle  \Phi f, k_{\lambda}^{\Theta}x\rangle}\nn\\
&=\overline{\langle  \Phi(\lambda) f(\lambda), x\rangle}\nn\\
&=\overline{\langle   f(\lambda), \Phi(\lambda)^{*} x\rangle}\nn\\
&=\overline{\langle   f, k_{\lambda}^{\Theta}\Phi(\lambda)^{*} x\rangle}\nn\\
&=\langle k_{\lambda}^{\Theta}\Phi(\lambda)^{*} x, f\rangle.\label{F1}
\end{align}
Therefore
$$(A_{\Phi}^{\Theta})^{*}k_{\lambda}^{\Theta}x= k_{\lambda}^{\Theta}\Phi(\lambda)^{*} x.$$
Since $\Phi,\Psi\in H^{\infty}(\oL)$, and  $$A_{\Phi}^{\Theta}=A_{\Psi}^{\Theta},$$ then by \cite[Corollary 6.4]{RT} there exist $k_{0}^{\Theta}\in M_{\Theta}$ and $X\in \oL$ such that 
$$\Psi=\Phi+k_{0}^{\Theta}X.$$
Observe that for each $\lambda\in \mathbb{D}$, $x\in E$, and $f\in K_{\Theta}$, 
\ben \langle k_{\lambda}^{\Theta}x, k_{0}^{\Theta}X f \rangle&=&\overline{\langle k_{0}^{\Theta}X f, k_{\lambda}^{\Theta}x\rangle}\nn\\
&=&\overline{\langle k_{0}^{\Theta}(\lambda)X f(\lambda), x\rangle}\nn\\
&=&\overline{\langle  f(\lambda), X^{*}(k_{0}^{\Theta}(\lambda))^{*}x\rangle}\nn\\
&=&\overline{\langle  f, k_{\lambda}^{\Theta}X^{*}(k_{0}^{\Theta}(\lambda))^{*}x\rangle}\nn\\
&=&\langle  k_{\lambda}^{\Theta}(k_{0}^{\Theta}(\lambda)X)^{*}x, f\rangle.\label{F2}\eeqn
From (\ref{F1}) and (\ref{F2}), we obtain
\begin{align*}
\langle (A_{\Psi}^{\Theta})^{*}k_{\lambda}^{\Theta}x, f \rangle
&=\langle k_{\lambda}^{\Theta}x, \Psi f \rangle\\
&=\langle k_{\lambda}^{\Theta}x, \Phi f \rangle+\langle k_{\lambda}^{\Theta}x, k_{0}^{\Theta}X f \rangle\\
&=\langle k_{\lambda}^{\Theta}\Phi(\lambda)^{*}x, f    \rangle+\langle  k_{\lambda}^{\Theta}(k_{0}^{\Theta}(\lambda)X)^{*}x, f\rangle\\
&=\langle k_{\lambda}^{\Theta}(\Phi(\lambda)^{*}x+k_{\lambda}^{\Theta}(k_{0}^{\Theta}(\lambda)X)^{*}x, f \rangle \\
&=\langle k_{\lambda}^{\Theta}(\Phi(\lambda)+k_{0}^{\Theta}(\lambda)X)^{*}x, f \rangle \\
&=\langle k_{\lambda}^{\Theta}\Psi(\lambda)^{*}x, f \rangle. 
\end{align*}
%{\color{blue} Now for any $f\in K_{\Theta}$, consider the inner product 
%\begin{align*}
%\langle (A_{\Psi}^{\Theta})^{*}k_{\lambda}^{\Theta}x, f \rangle
%&=\langle k_{\lambda}^{\Theta}x, \Psi f \rangle\\
%&=\langle k_{\lambda}^{\Theta}x, \Phi f \rangle+\langle k_{\lambda}^{\Theta}x, k_{0}^{\Theta}X f \rangle\\
%&=\overline{\langle \Phi f, k_{\lambda}^{\Theta}x  \rangle}+\overline{\langle k_{0}^{\Theta}X f, k_{\lambda}^{\Theta}x\rangle}\\
%&=\overline{\langle \Phi(\lambda) f(\lambda), x  \rangle}+\overline{\langle k_{0}^{\Theta}(\lambda)X f(\lambda), x\rangle}\\
%&=\overline{\langle f(\lambda),  \Phi(\lambda)^{*}x  \rangle}+\overline{\langle  f(\lambda), X^{*}(k_{0}^{\Theta}(\lambda))^{*}x\rangle}\\
%&=\overline{\langle f,  k_{\lambda}^{\Theta}\Phi(\lambda)^{*}x  \rangle}+\overline{\langle  f, k_{\lambda}^{\Theta}X^{*}(k_{0}^{\Theta}(\lambda))^{*}x\rangle}\\
%&=\langle k_{\lambda}^{\Theta}\Phi(\lambda)^{*}x, f    \rangle+\langle  k_{\lambda}^{\Theta}(k_{0}^{\Theta}(\lambda)X)^{*}x, f\rangle\\
%&=\langle k_{\lambda}^{\Theta}(\Phi(\lambda)^{*}x+k_{\lambda}^{\Theta}(k_{0}^{\Theta}(\lambda)X)^{*}x, f \rangle \\
%&=\langle k_{\lambda}^{\Theta}(\Phi(\lambda)+k_{0}^{\Theta}(\lambda)X)^{*}x, f \rangle \\
%&=\langle k_{\lambda}^{\Theta}\Psi(\lambda)^{*}x, f \rangle \\
%\end{align*}
Hence
$$(A_{\Psi}^{\Theta})^{*}k_{\lambda}^{\Theta}x=k_{\lambda}^{\Theta}\Psi(\lambda)^{*}x.$$ Since $$A_{\Phi}^{\Theta}=A_{\Psi}^{\Theta}$$ and the kernel map $x\to k_{\lambda}^{\Theta}x$ is injective, we get 
$$\Phi(\lambda)^{*}=\Psi(\lambda)^{*}, ~~\text{for all}~~\lambda\in \mathbb{D}.$$ Hence, $\Phi=\Psi$ as functions.

Next observe that, by the invariance of $ K_{\Theta}$ under $L_{\Phi}$,  for all $f\in K_{\Theta}$ we have $\Phi f\in K_{\Theta}$, which means that
$$0=\langle \Phi f, \Theta h \rangle=\langle f, \Phi^{*}\Theta h\rangle $$ for all $h\in H^{2}(E)$. It follows that $\Phi^{*}\Theta h\in \Theta H^{2}(E)$. By \cite[Lemma 2.2]{RT}, there exists $\Phi_{1}\in H^{\infty}(\mathcal{L}(E))$ such that $$\Phi^{*}\Theta=\Theta \Phi_{1}.$$
From this, using the identity $\Phi\Theta=\Theta\Phi$, we have
\begin{eqnarray} \Phi_{1}&=&\Theta^{*}\Phi^{*}\Theta\nn\\
&=&(\Phi\Theta)^{*}\Theta\nn\\
&=&(\Theta\Phi)^{*}\Theta\nn\\
&=&\Phi^{*}\Theta^{*}\Theta\nn\\
&=& \Phi^{*}\,.\label{Phi1=Phi*}\end{eqnarray}

\noindent Since $\Phi_{1}\in H^{\infty}(\mathcal{L}(E))$, $\Phi_{1}$ is analytic. Thus, from (\ref{Phi1=Phi*}) it follows that
$$\Phi^{*}\in H^{\infty}(\mathcal{L}(E)).$$
On the other hand, as observed above, $\Phi\in H^{\infty}(\mathcal{L}(E))$ is also analytic.  Hence, the only possibility is that $\Phi(\zeta)$ is constant.

The proof of the converse is straightforward. Hence it is left to the reader.
\end{proof}


%\begin{cor}
    %The unitary space $\h$ is $S_\Theta$-invariant if and only if $\Phi(\zeta)$ is constant.
%\end{cor}
%{\color{orange} I think here will be only "if" statement?}

%\begin{proof}
%The proof follows immediately from Theorems~\ref{inv} and \ref{Phi constant}.
%\end{proof}


\medskip



\medskip
%\begin{thm}
  % Let $\Phi\in L^{\infty}(\oL)$ such that $\|\Phi\|_{\infty}\leq1$ and $L_\Phi P_\Theta=P_\Theta L_\Phi$. Then, $A_{\Phi}$ is completely non-unitary on $\h^{\perp}=K_{\Theta_1}$.
%\end{thm}
 %\begin{proof}
    % Now suppose, toward contradiction, that $A_\Phi$ is not completely non-unitary on $K_{\Theta_1}$. Then there exists a nonzero closed subspace $M\subset K_{\Theta_1}$ which reduces $A_\Phi$ and on which $A_\Phi$ is unitary. But then $M$ is a unitary reducing subspace of $K_\Theta$. By the maximality of $\Theta_1 K_{\Theta_2}$ we have $M\subset \Theta_1 K_{\Theta_2}$. We get
     %$$M\subset \Theta_1 K_{\Theta_2}\cap K_{\Theta_1}.$$ But from the decomposition of $K_\Theta$ we have $$ \Theta_1 K_{\Theta_2}\cap K_{\Theta_1}=\{0\}$$ which follows that $M=\{0\}$. It is a contradiction. Hence, $A_\Phi$ is completely non-unitary on $K_{\Theta_1}$.
 %\end{proof}




\medskip
\section{\bf{Invertibility and solvability of matrix-valued asymmetric truncated Toeplitz operators}} \label{s4}
\subsection{Invertibility}
When the symbol $\varphi$ is the constant $1$, in \cite{HNP}, using Devinatz--Windom's theorem, %the authors  
Hru$\check{s}\check{c}$ev,  Nikol'skii,  Pavlov obtain the following link with the invertibility of a classical Toeplitz operator associated to a unimodular symbol. 

\begin{thm}\cite[Theorem 4]{HNP}\label{1}
Let $\theta_1,\theta_2$ be inner functions. The following statements are equivalent.
\begin{enumerate}
\item The operator $A_{1}^{\theta_{1},\theta_{2}}=P_{\theta_{2}}:K_{\theta_{1}}\to K_{\theta_{2}}$ is invertible.
\item The Toeplitz operator $T_{\overline{\theta}_{2}\theta_{1}}:H^{2}\to H^{2}$ is invertible.
\item $\mbox{dist}(\overline{\theta_1}\theta_2,H^\infty)<1$ and $\mbox{dist}(\overline{\theta_2}\theta_1,H^\infty)<1$.
\end{enumerate}
\end{thm}

\medskip
Recently in \cite[Theorem 5.6]{part3}, C\^{a}mara and Partington provide %get
 an interesting connection  between the invertibility of an asymmetric truncated Toeplitz operator and of  the Toeplitz operator with matrix symbol. Their theorem was proved for the domain $\mathbb{R}$, after establishing several lemmas in the scalar case. Below, we present a new %(algebraic)
 simplified (algebraic) proof of the matrix version for the unit circle $\mathbb{T}$ of their theorem. 
 We wish to emphasize that the result of C\^{a}mara and Partington is, however, more general for the scalar case.

\medskip

\begin{thm}\label{G} Let $\Theta_1$ and $\Theta_2$ be inner functions and 
 $\Phi\in L^{\infty}(\oL)$. Then $A_{\Phi}^{\Theta_{1},\Theta_{2}}$ is invertible if and only if the Toeplitz operator $T_{G}: H^2(E)\oplus H^2(E)\to H^2(E)\oplus H^2(E)$ is invertible,
where $G=\begin{pmatrix}
\Theta_{2}& \Phi\\
0&\Theta_{1}^{*}
\end{pmatrix}.
$
\end{thm}
\begin{proof}
Consider the equation generated by the operator $A_{\Phi}^{\Theta_{1},\Theta_{2}}\in \mathcal{T}(\Theta_{1},\Theta_{2})$, i.e.,
\begin{equation}\label{asymm}
A_{\Phi}^{\Theta_{1},\Theta_{2}}f_{\Theta_{1}}=P_{\Theta_{2}}(\Phi f_{\Theta_{1}})=g_{\Theta_{2}},
\end{equation}
where $f_{\Theta_{1}}\in K_{\Theta_{1}}$ and $g_{\Theta_{2}}\in K_{\Theta_{2}}$.
\medskip

The equation generated by Toeplitz operator $T_{G}$ with block matrix symbol is given by
 \begin{equation}\label{TG}
P_{+}\begin{pmatrix}
\Theta_{2}& \Phi\\
0&\Theta_{1}^{*}
\end{pmatrix}
\begin{pmatrix}
X_{1}\\
X_{2}
\end{pmatrix}
=\begin{pmatrix}
f_{1}\\
f_{2}
\end{pmatrix},
\end{equation}
where $X_{1}, X_{2}, f_{1}, f_{2}\in H^{2}(E)$.

Suppose $A_{\Phi}^{\Theta_{1},\Theta_{2}}:K_{\Theta_{1}}\to K_{\Theta_{2}}$ is invertible. We wish to prove that \eqref{TG} has unique solution. Performing block multiplication, this equation can be written as the system of two equations
%
 %We can write  \eqref{TG} as follows:
 \begin{equation}\label{first}
 P_{+}(\Theta_{2}X_{1})+P_{+}(\Phi X_{2})=f_{1}
 \end{equation}
 \begin{equation}\label{second}
 P_{+}(\Theta^{*}_{1}X_{2})=f_{2}.
 \end{equation}
 The general solution of \eqref{second} has the  form
 $$X_{2}=\psi+\Theta_{1}f_{2},$$
 where $\psi$ is an arbitrary function in %from
  $K_{\Theta_{1}}$. Substituting this into %Then
  \eqref{first} yields %has the form
 \begin{equation}\label{third}
 \Theta_{2}X_{1}+P_{+}(\Phi\psi)+P_{+}(\Phi\Theta_{1}f_{2})=f_{1}.
 \end{equation}
 Applying the projection operator $P_{\Theta_{2}}$ to both sides of \eqref{third}, we obtain
  \begin{equation}%\label{4th}
 P_{\Theta_{2}}P_{+}(\Phi\psi)+P_{\Theta_{2}}P_{+}(\Phi\Theta_{1}f_{2})=P_{\Theta_{2}}(f_{1}).\nn
 \end{equation}
Equivalently % It can be written as
 \begin{equation}\label{5th}
 A_{\Phi}^{\Theta_{1},\Theta_{2}} \psi=P_{\Theta_{2}}(f_{1}-P_{+}(\Phi\Theta_{1}f_{2})).
 \end{equation}
 By the invertibility of $A_{\Phi}^{\Theta_{1},\Theta_{2}}$,  %implies that
  equation \eqref{5th} has unique solution $\psi\in K_{\Theta_{1}}$, which in turn determines $X_2$. Since, $T_{\Theta_{2}}$ is isometric with range equal to $\Theta_2 H^2(E)$, by \eqref{third}, we have
  $$
 T_{\Theta_{2}}X_{1}=f_1-P_{+}(\Phi\psi)-P_{+}(\Phi\Theta_{1}f_{2})=h\in \Theta_2 H^2(E).
 $$
The equation
 \begin{equation}\label{X_1h}
     T_{\Theta_2}X_1=h
 \end{equation}
has at most one solution. 
Since $T_{\Theta_2}$  is injective and $h$ is in its range, there is unique $X_1\in H^2(E)$ satisfying \eqref{X_1h}. Explicitly, $$X_1=\Theta_2^{*} h\in H^2(E)$$
satisfying
 $$T_{\Theta_2}(X_1)=T_{\Theta_2}(\Theta_2^{*}h)=h.$$ 
 
Conversely, assume \eqref{TG} has unique solution for arbitrary $f_{1}, f_{2}\in H^{2}(E)$. We may express  \eqref{asymm} in the form
 \begin{equation}\label{atto}
 P_{+}(\Phi f_{\Theta_{1}})+\Theta_{2}X_{1}=g_{\Theta_{2}},
 \end{equation}
 where $X_{1}$ is some unknown function in %from
  $H^{2}(E)$. Along with \eqref{atto}, we consider an equation of the kind
 \begin{equation}\label{zero}
 P_{+}(\Theta_{1}^{*}f_{\Theta_{1}})=0.
 \end{equation}
 Rewriting \eqref{atto} and \eqref{zero} of the form:
 \begin{equation}\label{last}
P_{+}\begin{pmatrix}
\Theta_{2}& \Phi\\
0&\Theta_{1}^{*}
\end{pmatrix}
\begin{pmatrix}
X_{1}\\
f_{\Theta_{1}}
\end{pmatrix}
=\begin{pmatrix}
g_{\Theta_{2}}\\
0
\end{pmatrix},
\end{equation}
such that \eqref{last} has unique solution, we deduce that the operator $A_{\Phi}^{\Theta_{1},\Theta_{2}}$ is invertible. Hence, the required result follows.
\end{proof}
\subsection{Solvability}
     An operator $T\in \mathcal{L}(E)$ is said to be normally solvable if the range of $T$ is closed in $E$.
\medskip

%\ni {\color{red} I think we need a remark or a lemma explaining why a block triangular operator has closed range if and only if its block-diagonal entries have closed range, or give a reference.}

%    {\color{orange} 
    \begin{rem}\label{Blockrange}
    If $S, T\in \oL$ have closed range, then the block-triangular operator 
$$B=\begin{pmatrix}
    S&A\\
    0&T
\end{pmatrix}$$ also has closed range, provided that $A$ is bounded.
Moreover, range of 
$B$ can be described as
$$\text{Ran}~ (B)=\text{Ran}~(S)\oplus \text{Ran}~(T).$$
\end{rem}
%}

We now use Remark~\ref{Blockrange} to prove the following result.

\begin{prop} \label{prTG} Let $\Theta_1$ and $\Theta_2$ be inner functions and 
 $\Phi\in L^{\infty}(\oL)$. Then the operator $T_{G}$ is normally solvable.
\end{prop}
\begin{proof}
Observe that $T_{G}$ is the block-triangular Toeplitz operator
$$T_{G}=
\begin{pmatrix}
    T_{\Theta_2}& T_{\Phi}\\
    0& T_{\Theta^{*}_{1}}
\end{pmatrix}
$$
whose block diagonal entries are normally solvable. Indeed, since the symbol $\Theta_2$ is inner, the Toeplitz operator $T_{\Theta_2}$ is an isometry on $H^2(E)$, and hence has closed range. The operator $T_{\Theta^{*}_{1}}$ is a co-isometry, hence it has closed range also. 
By the boundedness of $T_\Phi$ and Remark \ref{Blockrange} it follows that $T_{G}$ has closed range and hence it is normally solvable.
\end{proof}
%{\color{blue}
In the following theorem we discuss the normal solvability of the asymmetric truncated Toeplitz operators.

%\ni {\color{orange} In all this section $\Phi\in L^{\infty}(\oL)$.}

\begin{thm} \label{ATTO} Let $\Theta_1$ and $\Theta_2$ be inner functions and 
 $\Phi\in L^\infty(\oL)$.  Then the operator $A_{\Phi}^{\Theta_1, \Theta_2}$ is normally solvable.
\end{thm}

\begin{proof}
To prove that $A_{\Phi}^{\Theta_1, \Theta_2}$ is normally solvable, we need to show that if $(k_n)\subset K_{\Theta_1}$ is a sequence of functions such that
    $$A_{\Phi}^{\Theta_1, \Theta_2}k_n\to h,$$
 then $h\in \text{Rang}(A_{\Phi}^{\Theta_1, \Theta_2})$.

Let $(k_n)$ be such a sequence. Then $(0, k_n)$ is in $H^2(E)\oplus H^2(E)$ and $T_{G}$ acts on 
$(0, k_n)$ as follows:
    $$T_{G}(0, k_n)=(P_{+}(\Phi k_n), P_{+}(\Theta^{*}_1 k_n)).$$    
Since $k_{n}\in K_{\Theta_1}$, we have
    $$P_{+}(\Theta^{*}_1 k_{n})=0.$$
    Hence,
    $$T_{G}(0, k_n)\to (h, 0)~~\text{in} ~~H^2(E)\oplus H^2(E).$$
    Since by Proposition \ref{prTG},  $T_{G}$ has a close range, the limit $(h, 0)$ must belong to $\text{Rang}(T_{G})$. So there exist $(f, g)\in H^2(E)\oplus H^2(E)$ such that $$T_{G}(f, g)=(h, 0).$$ Hence, we have
    $$(P_{+}(\Theta_2 f+\Phi g), P_{+}(\Theta^{*}_1 g))=(h, 0).$$
    From the equality of the second components, $P_{+}(\Theta^{*}_1 g)=0,$ it follows that  $g\in K_{\Theta_1}$. From the first components, we have
    $$P_{+}(\Theta_2 f+\Phi g)=h.$$
    Taking the projection on $K_{\Theta_2}$ we have
    $$P_{\Theta_2}(\Phi g)=h,$$ which is exactly
    $$A_{\Phi}^{\Theta_1, \Theta_2}(g)=h.$$ Hence,
    $\text{Rang}(A_{\Phi}^{\Theta_1, \Theta_2})$ is closed.
\end{proof}

%{\color{orange} I think the following corollary is not important now, because the operator $T_{G}$ is normally solvable? }
%\begin{cor}
%Let $\Theta_1$ and $\Theta_2$ be inner functions and $\Phi\in L^{\infty}(\oL)$. Then the operator $T_{G}$ is normally solvable if and only if $A_{\Phi}^{\Theta_1, \Theta_2}$ is normally solvable. 
%\end{cor}
%The proof follows from Proposition \ref{prTG} and Theorem \ref{ATTO}.

%Conversely, assume $A_{\Phi}^{\Theta_1, \Theta_2}$ has closed range. To prove that $T_{G}$ is normally solvable, we need to show that if $(f_n, g_n)\in H^2(E)\oplus H^2(E)$ such that 
%\ben T_G(f_n, g_n)\to (u, v)~~\text{in}~~H^2(E)\oplus H^2(E),\label{limTGfngn}\eeqn then $(u, v)\in \text{Rang}(T_{G})$.
%Let  $f_n$ and $g_n$ be chosen as above and use the identity
 %  $$T_{G}(f_n, g_n)=(P_{+}(\Theta_2 f_n+\Phi g_n), P_{+}(\Theta^{*}_1 g_n)).$$
%From this and (\ref{limTGfngn}), we obtain
 %   $$
  %  \begin{cases}
  % \ben &&P_{+}(\Theta^{*}_1 g_n)\to v,\label{limv}\\
  %  && %\quad 
  %  P_{+}(\Theta_2 f_n +\Phi g_n)\to u.\label{limu}\eeqn
%    \end{cases}
 %  $$
 %{\color{red} Above I percentaged out the command ``quad" since I don't see the reason for the misalignment.}
 
  %\ni Since $\Theta_1$ is inner,  $T_{\Theta^{*}_1}$ as a co-isometry, hence it has a closed range. From (\ref{limv}), %the equation $$P_{+}(\Theta^{*}_{1}g_n)\to v$$
   % we can find $g\in H^2(E)$ such that
   % \begin{equation}\label{v}
    %    P_{+}(\Theta^{*}_1g)=v.
   % \end{equation}


% Moreover, each of the functions $g_n$ can be written as
  % $$g_n=g+k_n,~~~\text{ for some } k_n\in K_{\Theta_1}.$$ Substituting this expression into (\ref{limu}) and rearranging the terms, we obtain
   %the first component equation we have
%  $$P_{+}(\Theta_2 f_n+\Phi g+\Phi k_n)\to u.$$
 %  After rearranging, we have
 % \ben P_{+}(\Theta_2 f_n+\Phi k_n)\to u-P_{+}(\Phi g).\label{lastlim}\eeqn
   %Taking the projection $P_{\Theta_2}$ onto $K_{\Theta_{2}}$ yields
  % $$P_{\Theta_2}(\Phi k_n)\to P_{\Theta_2}(u-P_{+}(\Phi g)).$$ 
  %Therefore
  % $$A_{\Phi}^{\Theta_1, \Theta_2}(k_n)\to h$$
  % where $h=P_{\Theta_2}(u-P_{+}(\Phi g))\in K_{\Theta_2}$.
%Since by assumption $A_{\Phi}^{\Theta_1, \Theta_2}$ has a closed range, %we have $h\in \text{Rang}(A_{\Phi}^{\Theta_1, \Theta_2})$. So 
  % there exists $k\in K_{\Theta_1}$ such that $$A_{\Phi}^{\Theta_1, \Theta_2}(k)=h.$$
  % {\color{blue}
  %From (\ref{lastlim}), we obtain the equation
  % $$P_{+}(\Theta_2 f+\Phi k)=u-P_{+}(\Phi g),$$
 %  or 
 %  \begin{equation}\label{u}
 %      P_{+}(\Theta_2 f+\Phi g+\Phi k)=u,
  % \end{equation}
  % From equations \eqref{v} and \eqref{u} we have 
 %  $(f, g+k)\in H^2(E)\oplus H^2(E)$ such that 
   %\begin{align*}
   % T_{G}(f, g+k)
   % &=(P_{+}(\Theta_2 f+\Phi g+\Phi k), P_{+}(\Theta_1^{*}g+\Theta_1^{*}k))\\
   % &=(P_{+}(\Theta_2 f+\Phi g+\Phi k), P_{+}(\Theta_1^{*}g))\\
   % &=(u, v), 
  % \end{align*}
% because $P_{+}(\Theta_1^{*}k)=0$ for all $k\in K_{\Theta_1}$. 
%Hence, $T_{G}$ is normally solvable. }








\begin{thebibliography}{80}

\bibitem{AM}  J. Angler, J. E. McCarthy, Pick Interpolation and Hilbert Function Spaces, Graduate Texts in Mathematics, Amer. Math. Soc. Vol 44, Providence, RI, 2002.

\bibitem{BT1} Ch. Benhida, D. Timotin, Functional models and finite dimensional perturbations of the shift, Integral Equ. Oper. Theory, {\bf 29} (1997), 187--196.

\bibitem{BT2} Ch. Benhida, D. Timotin, Operators invariant relative to a completely nonunitary contraction, Math. Z.
{\bf 299} (2021), no. 3-4, 1631--1649. 

\bibitem{BH} A. Brown, P.~R. Halmos, Algebraic properties of Toeplitz operators, J. Reine Angrew. Math.  {\bf 213} (1963-1964), 89--102.

%\bibitem{BL} J.A. Ball, A. Lubin, On a class of contractive perturbations of restricted shifts,  \emph{Pacific J. Math.} 63 (1976),  309--323.

\bibitem{part3}
M. C. C\^{a}mara, J.~R. Partington, Asymmetric truncated Toeplitz operators and Toeplitz operators with matrix symbols, J. Oper. Theory \textbf{77} (2017), no. 2, 455--479.
%\bibitem{BT2} Ch. Benhida, D. Timotin, Finite rank perturbations of contractions, \emph{Integral Equ. Oper. Theory} 36 (2000), 253--268.

\bibitem{ChFT}
I. Chalendar, E. Fricain, D. Timotin, A survey of some recent results on truncated Toeplitz operators, in Recent progress on operator theory and approximation in spaces of analytic functions, 59--77, Contemp. Math., {\bf 679}, Amer. Math. Soc., Providence, RI, 2016.

\bibitem{CFT} N. Chevrot, E. Fricain, D. Timotin, The characteristic function of a complex symmetric contraction,  Proc. Amer. Math. Soc. {\bf 135} (2007),  2877--2886.

\bibitem{GMR}
{S.~R. Garcia, J.~E. Mashreghi, W. Ross}, Introduction to Model Spaces and their Operators, Cambridge Studies in Advanced Mathematics, \textbf{148}, Cambridge University Press, Cambridge, 2016.

\bibitem{HNP}
{S. V. Hru$\check{s}\check{c}$ev, N.~K. Nikol'skii, B.~S. Pavlov}, Unconditional bases of exponentials and of reproducing kernels, Compl. Anal. Spec. Theory, \textbf{1979--80}, 214-337.

\bibitem{jl}
J. Jurasik, B. \L anucha, Asymmetric truncated Toeplitz operators on finite-dimensional spaces, Oper. Matrices \textbf{11} (2017), no. 1, 245--262.

\bibitem{RK}
R. Khan, The generalized Crofoot transform, Oper. Matrices {\bf15} (1), 225--237 (2021).

\bibitem{RAJ}
R. Khan, Y. Ameur, J. Khan, Characterizations of matrix valued asymmetric truncated Toeplitz operators, Rocky Mt. J. Math. {\bf 55} (4), 1049--1061 (2025).

\bibitem{RT}
{R. Khan, Dan. Timotin}, Matrix valued truncated Toeplitz operators: Basic Properties. Complex Anal. Oper. Theory {\bf 12} (2017), 997--1014.


%\bibitem{KL} E. Ko, and J. E. Lee, {\it On complex symmetric Toeplitz operator}, J. Math. Anal. Appl.  434(2016), 20-34.

%\bibitem{KKL} D. O. Kang, E. Ko, and J. E. Lee, {\it Remarks On complex symmetric operator matrices}, Linear Multilinear Algebra, 70 (2020), 1-11.
   % \bibitem{KLRem} E. Ko, and J. E. Lee, {\it Remarks On complex symmetric operator matrices}, Linear Multilinear Algebra, 67 (2019), 1198-1216.

%\bibitem{Pe} V.V. Peller, Hankel Operators and their Applications, Springer Verlag, New York, 2003.

\bibitem{VP}
V. I. Paulsen, M. Ragupathi, An Introduction to the Theory
of Reproducing Kernel Hilbert Spaces, Cambridge University Press, Cambridge, 2016.

\bibitem{Peller}
V. V. Peller, Hankel Operators and their Applications, Springer Verlag, New York, 2003.

 \bibitem{sar} D. Sarason, Algebraic properties of truncated Toeplitz operators,  Oper. Matrices {\bf 1}
(2007),  491--526.

\bibitem{sed}
N. A. Sedlock, Algebras of truncated Toeplitz operators, Oper. Matrices {\bf 5} (2011), 309--326.

%\bibitem{DT}
%D. Timotin, Redheffer product and characteristic functions, J. Math. Anal. App. (1995), {\bf196}, 823-840.

\bibitem{NF} B. Sz.-Nagy,  C. Foias, H. Bercovici, L. K\'erchy, Harmonic Analysis of Operators on Hilbert Space.  Revised and enlarged edition. Universitext. Springer, New York, 2010.

\end{thebibliography}













\end{document}






\noindent Rewayat Khan

\noindent  Istanbul Aydin University, Turkey.
\noindent e-mail: rewayat.khan@gmail.com

\vspace{5mm}

\noindent Flavia Colonna

\noindent   Department of Mathematical Sciences, George Mason University, Fairfax, Virginia, U.S.A.

\noindent e-mail: fcolonna@gmu.edu






























